\section{Limitations}
\label{sec:limitations}

The first limitation concerns time variation and information availability. The RMT and network features used in the initial volatility-modelling exercise are estimated from the full synchronized sample and then broadcast across dates. This construction uses information unavailable at an historical forecast origin; consequently, the results for feature sets containing these descriptors cannot be interpreted as strict out-of-sample or real-time forecast performance. They are useful for identifying an ex-post structural association, but a valid forecasting test must reconstruct RMT and network features recursively, using only observations available at each origin, and re-select models within an expanding or rolling evaluation design.

The second limitation concerns the target set. We use PETR4, VALE3 and BBDC4 in the forecasting experiment. These assets are liquid and economically important, but they do not cover the full B3 cross-section. Future work should extend the forecasting exercise to a larger set of assets and evaluate whether network features help more for some sectors than others.

The third limitation concerns data quality. The core historical universe is sufficiently clean for the current analysis, but broader or more liquid modern universes may contain extreme adjusted-return events that require additional validation. These events may reflect corporate actions, ticker transitions, units or adjustment-factor issues. We would flag, inspect and either correct or handle them through robustness checks before using broader universes for final network claims.

The fourth limitation concerns sector classification and issuer representation. The macro-sector and subsector map is an initial classification. It is adequate for testing whether within-sector correlations exceed between-sector correlations, but it should be reviewed ticker by ticker before submission. We provide an issuer-collapsed RMT robustness check that removes duplicate share classes; nevertheless, the deterministic representative rule is not a substitute for a fully specified liquidity, free-float or economic-exposure rule. This remains especially important for holding companies, units and firms with mixed exposure across sectors.

Finally, the forecasting analysis is limited to three target assets, static network features and simple neural architectures. Feature importance in tree models is useful but does not establish causal mechanisms. PMFG features show nonzero but modest importance, and the magnitude is small compared with lagged-volatility predictors. Future work should include rolling network features, a broader target asset set, richer temporal neural architectures and graph neural networks.
